% !TeX root = ../report_jouleecosystem.tex
\chapter{The Bridge: from measurements to decisions}

While the three repositories are independent, this short chapter shows how they work together end-to-end and how data flows between them.

\section{The pipeline}

\begin{enumerate}
    \item \textbf{Measure} (JouleQuest): A testing campaign on a board generates raw data (CSV/JSON). We process this into a summary, and finally into an energy lookup table for that specific board.
    \item \textbf{Estimate} (JouleGrad): We load the lookup table into the \texttt{EnergyEstimator}. This tool provides differentiable energy estimates for the layer shapes we measured. By default, it refuses to guess for unmeasured shapes.
    \item \textbf{Search} (JouleNAS): The estimator is added to the training loss as $\lambda_E \widehat{E}$. The optimizer then prunes channels that cost more energy than they provide in accuracy, using the real measured costs of the target board.
\end{enumerate}

The key strength of this pipeline is that no numbers are guessed. The energy values driving the search come from real, physical measurements on the actual hardware, rather than theoretical FLOP counts or datasheet averages.

\section{From masks to physical deployment: validating with JouleQuest}
\label{sec:joulenas_to_joulequest}

So far, JouleNAS has evaluated pruned models by learning channel masks. While mask training is convenient and differentiable, it is still an emulation: the underlying PyTorch network still executes full-size dense convolutions and simply multiplies pruned channels by zero. To see real hardware energy savings and measure them on our NVIDIA Jetson board, we need to take the next step: \emph{deploying a physically smaller model}.

\subsection{Selecting the target models: \texorpdfstring{$\lambda_E = 0.005$}{lambda = 0.005} discrete}

To validate our methodology on physical hardware, we selected the CIFAR-10 ResNet-18 checkpoints obtained with $\lambda_E = 0.005$ in discrete mode for both benchmarked devices:
\begin{itemize}
    \item \textbf{NVIDIA Jetson AGX Orin}: Preserves $92.10\%$ test accuracy (compared to the $92.30\%$ dense baseline, a loss of only $0.20\%$) with a predicted energy fraction of $58.27\%$ (a $41.73\%$ energy saving).
    \item \textbf{Raspberry Pi 5}: Preserves $92.27\%$ test accuracy (compared to the $92.31\%$ dense baseline, losing only $0.04\%$) with an estimated energy fraction of $22.54\%$ (a $77.46\%$ predicted energy saving), pruning down to $60.40\%$ active channels.
\end{itemize}
Both checkpoints offer substantial predicted energy reductions well above physical measurement noise and environmental fluctuations.

\subsection{Reconstructing the hardware-specific pruned architectures}

Using the per-layer pruning logs exported by JouleNAS, we reconstructed the materialized models inside JouleQuest:
\begin{itemize}
    \item \texttt{PrunedOrinCifarResNet18} (NVIDIA Jetson AGX Orin):
    \begin{itemize}
        \item \textbf{Stem}: $3 \to 64$.
        \item \textbf{Stage 1}: Block 1.0 is dense ($64 \to 64 \to 64$), Block 1.1 prunes its mid-convolution to 58 ($64 \to 58 \to 64$).
        \item \textbf{Stage 2}: Block 2.0 uses $64 \to 128 \to 127$ (shortcut $64 \to 127$); Block 2.1 uses $127 \to 128 \to 127$.
        \item \textbf{Stage 3}: Block 3.0 uses $127 \to 253 \to 249$ (shortcut $127 \to 249$); Block 3.1 uses $249 \to 192 \to 236$.
        \item \textbf{Stage 4}: Block 4.0 uses $236 \to 381 \to 222$ (shortcut $236 \to 222$); Block 4.1 prunes its internal layer down to 16 channels ($222 \to 16 \to 297$).
        \item \textbf{Classifier}: Linear $297 \to 10$.
    \end{itemize}
    \item \texttt{PrunedPi5CifarResNet18} (Raspberry Pi 5):
    \begin{itemize}
        \item \textbf{Stem}: $3 \to 64$.
        \item \textbf{Stage 1}: Both Block 1.0 and Block 1.1 remain dense at 64 channels ($64 \to 64 \to 64$).
        \item \textbf{Stage 2}: Block 2.0 uses $64 \to 127 \to 128$ (shortcut $64 \to 128$); Block 2.1 uses $128 \to 127 \to 128$.
        \item \textbf{Stage 3}: Block 3.0 uses $128 \to 248 \to 231$ (shortcut $128 \to 231$); Block 3.1 uses $231 \to 139 \to 238$.
        \item \textbf{Stage 4}: Block 4.0 uses $238 \to 127 \to 248$ (shortcut $238 \to 248$); Block 4.1 uses $248 \to 64 \to 258$.
        \item \textbf{Classifier}: Linear $258 \to 10$.
    \end{itemize}
\end{itemize}

\subsection{Handling residual shortcuts and BatchNorm}

Reconstructing a residual architecture with variable layer sizes requires two practical adjustments:
\begin{enumerate}
    \item \textbf{Residual shortcut alignment}: In standard ResNet, identity shortcuts simply pass inputs forward when stride is 1. But when a block's input channels do not equal its output channels (such as Block 3.1 and Block 4.1 on both platforms), an element-wise addition would encounter mismatched shapes. In our reconstructed blocks, the shortcut automatically applies a $1 \times 1$ convolution to match channel dimensions seamlessly.
    \item \textbf{BatchNorm adjustments}: Each BatchNorm layer is sized to match the exact active channels of its preceding convolution. While the JouleNAS energy estimate did not include BatchNorm or residual additions in its 21-term lookup, the convolutional and linear layers account for the vast majority of computation.
\end{enumerate}

\subsection{Automated testing and validation in JouleQuest}

To make physical testing straightforward and repeatable, we provided dedicated validation suites in JouleQuest's campaign launcher (\texttt{run\_measurement\_campaign.sh}):
\begin{itemize}
    \item For Jetson AGX Orin:
\begin{lstlisting}
./run_measurement_campaign.sh --board agx_orin --suite pruned_validation_orin
\end{lstlisting}
    \item For Raspberry Pi 5:
\begin{lstlisting}
CPU_THREADS=4 BACKEND=cpu SLEEP_TIME=10 ./run_measurement_campaign.sh --board pi5_4threads --suite pruned_validation_pi5
\end{lstlisting}
\end{itemize}

Each suite schedules exactly two targeted experiments:
\begin{enumerate}
    \item The dense CIFAR-10 ResNet-18 baseline (\texttt{ResNet18\_32\_10.pt});
    \item The hardware-specific materialized pruned model (\texttt{PrunedOrinResNet18\_32\_10.pt} or \texttt{PrunedPi5ResNet18\_32\_10.pt}).
\end{enumerate}

Additionally, the comparison tool \texttt{validate\_pruned\_measurements.py} automatically matches the physical INA226 measurements against the JouleNAS predictions, quantifying the measured energy reduction and the composition discrepancy.

Physical testing on the Raspberry~Pi~5 revealed two key operational parameters for CPU-based embedded devices:
\begin{itemize}
    \item \textbf{Inter-burst cooldown}: With a brief 3-second sleep, residual thermal buildup and dynamic CPU frequency scaling distorted dense baseline calibration. Extending the sleep period to 10~seconds allowed the CPU to cool completely between bursts.
    \item \textbf{CPU thread allocation}: Because the original layer lookup table on the Raspberry~Pi~5 was profiled utilizing the full quad-core Cortex-A76 processor, we explicitly enforced 4 CPU threads (\texttt{CPU\_THREADS=4}) on the measurement runner to maintain strict hardware parity.
\end{itemize}
Under these stabilized, 4-thread conditions:
\begin{itemize}
    \item The dense ResNet-18 baseline consumed $28.38$\,mJ per inference (compared to the table-predicted $35.15$\,mJ, a $19.3\%$ composition gap).
    \item The materialized \texttt{PrunedPi5ResNet18} consumed $11.33$\,mJ per inference (predicted $7.93$\,mJ).
    \item Physical measurements showed a \textbf{60.08\% real-world energy reduction} on physical hardware (compared to the $77.46\%$ theoretical prediction), with only a negligible $0.04\%$ loss in test accuracy ($92.27\%$ vs $92.31\%$).
\end{itemize}

\subsection{Physical hardware validation on NVIDIA Jetson AGX Orin (batch size 32)}
\label{sec:bridge_agxorin_bs32_validation}

\noindent\textit{[Note: This subsection was written by the AI assistant to document and analyze the experimental validation results on NVIDIA Jetson AGX Orin at batch size 32.]}
\vspace{0.5em}

To close the loop on the embedded GPU target and explore batch scaling behavior on physical silicon, we deployed the materialized \texttt{PrunedOrinBs32ResNet18} (derived from the optimal $\lambda_E = 1.20$ checkpoint) alongside the dense baseline on the physical NVIDIA Jetson AGX Orin at batch size 32. Measurements were orchestrated using JouleQuest's dedicated validation pipeline:
\begin{lstlisting}
./validate_agx_orin_bs32.sh
# or equivalently via the campaign orchestrator:
NETWORK_BATCH_SIZE=32 ./run_measurement_campaign.sh --board agx_orin_bs32 --suite pruned_validation_orin_bs32
\end{lstlisting}

Physical power and energy traces recorded via the INA226 sensor yielded the following results:
\begin{itemize}
    \item \textbf{Dense baseline fidelity}: The measured dense ResNet-18 consumed \textbf{$1.368$\,mJ per sample} (mean burst power: $8.07$\,W, burst duration: $16.8$\,ms). This matches the JouleNAS lookup table prediction of $1.372$\,mJ to within \textbf{$-0.3\%$} (a discrepancy of only $-0.004$\,mJ). This near-perfect alignment provides an encouraging hint to consider, suggesting that at batch size 32, launch overheads are amortized and isolated layer lookups compose much closer to the full architecture execution.
    \item \textbf{Materialized pruned model performance}: The physical model (\texttt{PrunedOrin\-Bs32\-ResNet18}) consumed \textbf{$1.095$\,mJ per sample} (mean burst power: $5.79$\,W, burst duration: $15.2$\,ms).
    \item \textbf{Real-world energy and power savings}: Physical measurements demonstrated a confirmed \textbf{$19.95\%$ real-world energy reduction} (saving $0.273$\,mJ per inference across all 32 samples) and a \textbf{$28.25\%$ reduction in average power dissipation} (dropping from $8.07$\,W to $5.79$\,W), while comfortably preserving $89.26\%$ classification accuracy (an accuracy delta of only $-2.60\%$).
\end{itemize}

\subsubsection{Analysis of the physical vs. predicted savings gap}
While JouleNAS predicted a theoretical energy reduction of $34.58\%$ ($0.898$\,mJ), physical hardware measurements yielded $19.95\%$ ($1.095$\,mJ), reflecting a composition discrepancy of $+22.0\%$ on the pruned architecture. Several hardware factors explain this residual difference on embedded GPU architectures:
\begin{enumerate}
    \item \textbf{Unaligned warp execution and Tensor Core tiling}: Deep structured pruning reduces channels to non-standard widths (such as $58, 127, 249, 236, 222$). Because modern NVIDIA Ampere architectures optimize GEMM routines for channel dimensions that are multiples of 32 or 64, irregular filter counts cause partial warp execution and lower compute efficiency per SM.
    \item \textbf{Residual additions and BatchNorm}: Reconstructing the residual topology requires additional $1 \times 1$ projection convolutions on shortcuts where channel dimensions mismatch, as well as BatchNorm layers whose costs were not accounted for in the isolated 21-term convolutional lookup.
    \item \textbf{Static platform power floor}: Even during active computation, the Jetson AGX Orin baseline SoC components (memory controllers, CPU host interface, and PCIe peripherals) maintain a baseline idle power floor of $\approx 2.1$\,W, which dampens proportional energy reductions when execution time drops.
\end{enumerate}
Despite these low-level GPU realities, a physical energy reduction of $20\%$ accompanied by a $2.28$\,W drop in power dissipation provides an encouraging practical hint that JouleNAS can produce hardware-pruned networks that save substantial energy on real edge devices.

\section{Where the chain is weakest}

Two weak points need attention:

\begin{itemize}
    \item \textbf{Composition}: We build the tables from isolated layer measurements, but the search uses them in full networks. This leads to a gap between the sum of the parts and the true total (Section \ref{sec:sum-of-parts}). The gap is small on Jetson boards but reaches up to about $25\%$ on the Raspberry~Pi~5.
    \item \textbf{Coverage}: The estimator only knows about the shapes we actually measured. If the search needs an unmeasured shape, the system refuses to guess by default. This is a deliberate safety feature so we don't trust unmeasured numbers.
\end{itemize}

Keeping these points in mind, the pipeline is ready to use and, as the next chapter explains, ready to be improved.